\chapter{Overview}

\section{Problem Statement}

Operational procedures are the connective tissue of production engineering. When a service degrades,
an incident fires, or a change must be applied under pressure, engineers rely on runbooks---step-by-step
guides that encode institutional knowledge about how to respond. The research literature is unambiguous:
\textit{unstructured runbooks fail.} Static documents drift from reality, skip governance checks, and
produce no audit evidence. Teams that operate at scale cannot afford a response that depends entirely on
whether the on-call engineer remembers to follow the wiki~\cite{ibm-redbook-runbooks,ijsret-runbook-engineering}.

The core operational problem is threefold. First, runbooks must be \textbf{executable}: a document that
cannot be run is a suggestion, not a procedure. Second, runbooks must be \textbf{governed}: every action
that touches infrastructure must be bounded by policy---command allowlists, environment variable
controls, output redaction, and role-based approval gates. Third, runbooks must be \textbf{traceable}:
every execution must produce an immutable, auditable record that satisfies SOC~2, ISO~27001, and HIPAA
audit requirements~\cite{soc2-aicpa,iso27001,hipaa-security-rule} without manual effort. Without all three properties, organizations face audit
failures, runbook drift, and uncontrolled blast radius during incidents.

Gert is designed to solve exactly this problem: a governed, executable, debuggable runbook engine that
produces full traceability and evidence by construction.

\section{Why v2: What v1 Got Right, and What Constrained It}

Gert v1 validated the core thesis. It delivered governed execution with command allowlists, environment
variable blocking, output redaction, and per-step approval gates. It captured evidence---text,
checklists, attachments with SHA256 hashes---and wrote an append-only JSONL trace for every run. It
provided deterministic replay via recorded scenarios, a full terminal UI (Bubble Tea), a VS~Code
extension over JSON-RPC~2.0, interactive debugging with a step-through REPL, and a \texttt{gert~test}
scenario runner. These are production-grade features that must be preserved in v2.

However, v1 accumulated architectural constraints that limit its evolution. The serve layer
(\texttt{gert serve}) is tightly coupled to the runtime, making it difficult to add new adapters without
modifying core logic. Schema versioning is brittle: \texttt{runbook/v0} and \texttt{runbook/v1} diverge
with no formal compatibility contract. The tool and provider subsystems are functional but have no stable
extension contract, making third-party tool development fragile. There is no saga or compensation pattern
for rollback on failure, no parallel fan-out execution, no policy-as-code integration, and no
OpenTelemetry support for distributed tracing---all of which are now industry-standard in comparable
workflow orchestration systems~\cite{temporal-docs,argo-workflows,prefect-docs}.

These constraints are architectural, not cosmetic. A clean-break v2 is warranted.

\section{Scope}

\textbf{Gert v2 is:} a local-first, governed runbook execution engine with a stable component
architecture, versioned extension contracts, a first-class governance and policy layer, full traceability
by construction, and defined adapter interfaces for TUI, VS~Code, and web clients. It ships as a single
Go binary with an optional out-of-process extension runtime.

\textbf{Gert v2 is not:} a managed cloud execution service, a general-purpose workflow orchestrator, a
CI/CD pipeline system, a SOAR platform, or a replacement for systems like Temporal, Argo, or PagerDuty
Runbook Automation~\cite{temporal-docs,argo-workflows,pagerduty-automation}. It does not provide a central policy server, a shared execution registry, or
multi-tenant execution isolation.

\section{Design Priorities}

\begin{itemize}
  \item \textbf{A core runtime isolated from presentation adapters.} The execution engine---parser,
    planner, state machine, governance enforcer, trace writer---must have no dependency on any
    presentation layer. TUI, VS~Code, and web clients attach via a stable event subscription API;
    they observe execution but never own it. This isolation is the prerequisite for every other
    architectural goal.

  \item \textbf{Explicit and versioned contracts for schema, tools, and events.} Every interface that
    crosses a component boundary must be named, versioned, and documented. The runbook schema carries
    an \texttt{apiVersion} field with formal compatibility guarantees. Tool definitions carry a version.
    The runtime event stream has a defined schema with ordered delivery guarantees. Breaking a contract
    requires a version increment, not a silent edit.

  \item \textbf{A secure extension runtime using out-of-process plugins.} Extensions run in isolated
    processes and are granted capabilities explicitly. No extension receives ambient authority. The
    capability model---what an extension may register, observe, or invoke---is a named, versioned list
    that the host enforces at admission time. This pattern follows the principle of least privilege
    endorsed by NIST SP~800-53~\cite{nist-sp800-53} and applied in production by systems such as Kubernetes admission
    controllers~\cite{kubernetes-docs} and OPA Gatekeeper~\cite{opa-gatekeeper}.
\end{itemize}

\section{Document Roadmap}

This document is organized as follows. \S\ref{ch:goals} defines measurable goals and explicit
non-goals, establishing scope boundaries for implementors. \S\ref{ch:architecture} describes the five
primary components and the dependency rules between them. \S\ref{ch:schema} specifies the v2 runbook
schema, including versioning strategy, field inventory, and migration path from v1. \S\ref{ch:extension}
defines the extension runtime: capability taxonomy, handshake protocol, and trust model.
\S\ref{ch:tools} specifies the tool runtime: discovery, invocation, transport contracts, and governance
integration. \S\ref{ch:events} defines the runtime event model: event schema, ordering guarantees,
delivery semantics, and replay contract. \S\ref{ch:security} specifies the security and trust model:
threat model, policy enforcement, supply chain controls, and audit requirements. \S\ref{ch:testing}
defines the testing and acceptance strategy: contract test categories, coverage requirements, and
acceptance criteria. \S\ref{ch:questions} documents open architectural questions with options and
resolution criteria.

\section{Success Definition}

Gert v2 is complete when: (1) a v1 runbook can be migrated to v2 schema via \texttt{gert migrate} and
executed without manual fixup; (2) the TUI, VS~Code extension, and a third-party adapter can all be
built against the published event and API contracts without importing internal packages; (3) the
governance layer---allowlists, denylists, env blocking, redaction, approval gates---is enforced for
every execution mode including extension-invoked runs; (4) every execution produces an append-only,
crash-safe JSONL trace that is structurally valid against the v2 trace schema; and (5) \texttt{gert
test} passes a scenario suite covering all step types, governance rules, and evidence capture cases.
